\section{System Model and Problem Formulation}
\label{sec:system}

\subsection{Network and Slicing Model}

We consider an operator infrastructure divided into a set of
\emph{reservable pools} --- a radio access network partitioned into
sectors, or a slicing-enabled transport/compute fabric partitioned
into slices. Time is slotted at the operational control granularity
$\delta$; throughout this paper $\delta = 15$ minutes, the standard
PM-counter granularity of 3GPP performance management and of our
dataset. Each pool $i$ has a physical capacity $C_i$: the maximum
demand it can serve in one slot. The aggregate demand offered to pool
$i$ in slot $t$ is $\Dt^{(i)} \ge 0$, a realization of a stochastic
process with strong diurnal and weekly seasonality, heavy-tailed
within-slot bursts, and occasional regime shifts
(mass events, flash crowds, gradual traffic evolution).

At each slot $t$ the controller commits a \emph{reservation}
$\Rt^{(i)} \in [0, C_i]$ for pool $i$ covering slot $t+1$. The
reservation is a hard commitment in the operational sense: demand up
to $\Rt$ is served with contracted quality; demand beyond $\Rt$
violates the SLA (blocked, degraded, or served from best-effort
overflow at a quality penalty). Reservations may be revised every
slot, but large swings are themselves costly --- reconfiguration
signaling, spectrum map churn, and upstream commitment churn --- which
we model with a switching penalty implicitly smoothed by the update
mechanism of Section~\ref{sec:stage4}.

\subsection{Closed-Loop Decision Process}

The controller observes, at the end of slot $t$: the demand history
$\{\D_{t-k}\}_{k=0}^{K-1}$ (a window of $K$ past slots; $K=96$, i.e.,
24 hours), the realized outcome of past reservations
$\mathbf{1}[\D_{t} > R_{t}]$, and auxiliary feedback (recent
violation rate, demand level relative to training statistics). It
does \emph{not} observe $\D_{t+1}, \D_{t+2}, \dots$. This is the
defining constraint of the problem and, as
Section~\ref{sec:leakage-audit} details, the constraint most often
violated inadvertently in experimental harnesses. Formally, the
information set available when choosing $R_{t+1}$ is
\begin{equation}
\mathcal{F}_t \;=\; \sigma\bigl(\D_{t-K+1..t},\;
\{R_{s}\}_{s\le t},\; \{\mathbf{1}[\D_s > R_s]\}_{s \le t}\bigr),
\label{eq:filtration}
\end{equation}
and any admissible policy must be $\mathcal{F}_t$-measurable. All
components of UCRA --- the forecaster, the transform, the risk state,
the replay buffer --- are functions of $\mathcal{F}_t$ only; the unit
tests of Section~\ref{sec:setup} enforce this mechanically.

\subsection{Reservation Metrics}

For a pool with capacity $C$ and a reservation trajectory
$\{R_t\}_{t=1}^{T}$ against realized demand $\{\D_t\}_{t=1}^{T}$,
we use four closed-loop metrics.

\emph{Violation rate} (risk exposure):
\begin{equation}
V \;=\; \frac{1}{T}\sum_{t=1}^{T} \mathbf{1}[\D_t > R_t].
\label{eq:violation}
\end{equation}
This is the SLA-failure frequency --- the quantity the operator is
contractually penalized for.

\emph{Utilization} (efficiency of the committed capacity):
\begin{equation}
U \;=\; \frac{1}{T}\sum_{t=1}^{T} \min\!\left(\frac{\D_t}{R_t},\,1\right).
\label{eq:utilization}
\end{equation}
Per-slot utilization is capped at 1: over-demand does not inflate
utilization; it appears as violations instead. A policy can therefore
exhibit both high utilization and high violation rate
(median-forecast policies do exactly this on our data).

\emph{Waste} (idle committed capacity):
\begin{equation}
W \;=\; \frac{1}{T\,C}\sum_{t=1}^{T} \max(R_t - \D_t,\, 0),
\label{eq:waste}
\end{equation}
the fraction of physical capacity that was reserved, went unused, and
could have been monetized by other tenants.

\emph{Reservation error} (tracking):
\begin{equation}
E \;=\; \frac{1}{T\,C}\sum_{t=1}^{T} |R_t - \D_t|,
\label{eq:reserr}
\end{equation}
the mean absolute gap, normalized. We emphasize --- and our results
demonstrate --- that $E$ alone is \emph{not} a meaningful objective:
a policy that tracks demand with the median minimizes $E$ while
violating half the time. $V$, $U$, and $W$ are the operational
objectives; $E$ is diagnostic.

\subsection{Objective}

The operator's problem is to choose an
$\mathcal{F}_t$-measurable reservation policy $\pi$ to
\begin{equation}
\min_{\pi}\; \mathbb{E}\bigl[W(\pi)\bigr]
\quad \text{subject to} \quad
V(\pi) \;\le\; \epsilon,
\label{eq:objective}
\end{equation}
i.e., minimize waste subject to a violation budget $\epsilon$
(typically $\epsilon \le 0.01$ for service-grade slices, and
$\epsilon = 0$ is desirable for URLLC-grade commitments). This is a
chance-constrained sequential decision problem with unknown dynamics;
exact solutions are intractable without a full demand model, and
learned policies raise the auditability concerns discussed earlier.
UCRA's design stance is to solve the problem \emph{implicitly}:
a calibrated quantile forecast makes $V$ approximately
$(1-\tau)$-controllable by construction; $\kappa$ buys the remaining
headroom that forecast errors consume; and the closed loop adapts the
forecast when the world drifts. The frontier
$(V(\kappa), U(\kappa))$ is then an empirical property of the
 deployed system, which we measure in full.

\subsection{Threat Model for Evaluation}

A reservation policy can cheat in three ways: (i) by using future
demand anywhere in its decision path (forecast leakage); (ii) by
using the \emph{current} slot's realized demand to size its own
reservation (outcome leakage --- a policy cannot know what it must
protect before the protection is committed); and (iii) by adapting on
labels that are not yet observable (label-availability leakage).
Category (iii) is the subtle one that motivates
Section~\ref{sec:stage5}: an online-updated model is entitled to learn
from history, but only from the part of history that has already
happened. Our evaluation protocol is designed to make all three
categories detectable and, in our final pipeline, impossible.

\diagramplaceholder{4.2cm}{ucra-architecture}
{Suggested content: the five UCRA stages arranged as a closed loop
over the network --- Stage 1 (Quantile LSTM) feeding Stage 2
(uncertainty-to-reservation transform $\Phi$), feeding Stage 3
(risk-adaptive allocation), feeding the network/slices; reservation
and demand feedback returning through Stage 4 (EWMA update and
release hysteresis) and Stage 5 (drift monitor, replay buffer,
fine-tuning). Mark the information boundary ``observable at time $t$''
around the feedback path.}
